\subsection{线性表示的权}

在本节中 k 表示域 R 或 C 之一。

设 V 是 k 上的有限维向量空间，\(\rho: G \to \mathbf{GL}(V)\) 是连通紧李群 G 在 V 上的连续（从而是实解析的，第 III 章，§8，第 1 节，定理 1）表示。如下定义一个复向量空间 \(\tilde{V}\) 与连续表示 \(\tilde{\rho}: G \to \mathbf{GL}(\tilde{V})\) ：若 k = C，令 \(\tilde{V} = V\) ，\(\tilde{\rho} = \rho\) ；若 k = R，令 \(\tilde{V} = V_{(C)}\) ，并令 \(\tilde{\rho}\) 为 \(\rho\) 与典则同态 \(\mathbf{GL}(V) \to \mathbf{GL}(\tilde{V})\) 的复合。

对一切 \(\lambda \in \mathrm{X}(\mathrm{G})\) ，用 \(\tilde{\mathrm{V}}_{\lambda}(\mathrm{G})\) 表示 \(\tilde{\mathrm{V}}\) 中由诸 \(v \in \tilde{\mathrm{V}}\) （满足 \(\tilde{\rho}(g)v = g^{\lambda}v\) 对一切 \(g \in \mathrm{G}\) ）组成的向量子空间（参见第 VII 章，§1，第 1 节）。由同前命题 3，诸 \(\tilde{\mathrm{V}}_{\lambda}(\mathrm{G})\) （\(\lambda\) 取遍 \(\mathrm{X}(\mathrm{G})\) ）之和是直和。此外：

引理 1. 若 G 交换，则 \(\tilde{V}\) 是诸 \(\tilde{\mathrm{V}}_{\lambda}(\mathrm{G})\) （\(\lambda \in \mathrm{X}(\mathrm{G})\) ）的直和。

由于 \(\rho\) 半单（§1 第 1 节），只需在 \(\rho\) 单的情形证明引理。这时，交换子 \(Z\) （\(\rho(G)\) 在 \(\mathrm{End}(\tilde{V})\) 中的）约化为位似（《代数学》，第 VIII 章，§3，第 2 节，定理 1）；于是，同态 \(\tilde{\rho}\) 的像含于子群 \(\mathbf{C}^*.1_V\) （\(\mathrm{GL}(\tilde{V})\) 的），且存在 \(\lambda \in X(G)\) 使得 \(\tilde{V} = \tilde{V}_{\lambda}(G)\) 。

定义 1. G 的表示 \(\rho\) 相对于 G 的极大环面 T 的权，是元素 \(\lambda\) （属于 X(T)，且使 \(\tilde{\mathrm{V}}_{\lambda}(\mathrm{T}) \neq 0\) ）。

用 \(\mathrm{P}(\rho,\mathrm{T})\) （若 T 的选择不致混淆，也简记为 \(\mathrm{P}(\rho)\) ）表示 \(\rho\) 相对于 T 的权所成的集。由引理 1，

\[
\tilde {\mathrm{V}} = \bigoplus_ {\lambda \in \mathrm{P} (\rho , \mathrm{T})} \tilde {\mathrm{V}} _ {\lambda} (\mathrm{T}).\tag{7}
\]

设 \(T'\) 是 G 的另一个极大环面，g 是满足 \((\operatorname{Int} g)T = T'\) 的 G 的元素（§2 第 2 节，定理 2）。对一切 \(\lambda \in X(T)\) ，

\[
\tilde {\rho} (g) (\tilde {\mathrm{V}} _ {\lambda} (\mathrm{T})) = \tilde {\mathrm{V}} _ {\lambda^ {\prime}} (\mathrm{T} ^ {\prime}), \quad \text { where } \lambda^ {\prime} = \mathrm{X} (\operatorname{Int} g ^ {- 1}) (\lambda).\tag{8}
\]

从而，

\[
\mathrm{X(Int} g) (\mathrm{P} (\rho , \mathrm{T} ^ {\prime})) = \mathrm{P} (\rho , \mathrm{T}).\tag{9}
\]

Weyl 群 \(\mathrm{W} = \mathrm{W}_{\mathrm{G}}(\mathrm{T})\) 左作用在 \(\mathbf{Z}\) -模 \(\mathrm{X(T)}\) 上，作用由 \(w\mapsto \mathrm{X}(w^{-1})\) 给出；于是，对 \(t\in \mathrm{T},\lambda \in \mathrm{X(T)},w\in \mathrm{W}\) ，有 \(t^{w\lambda} = (w^{-1}(t))^{\lambda}\) 。

命题 5. 集 \(\mathrm{P}(\rho, \mathrm{T})\) 在 Weyl 群 W 的作用下稳定。设 \(n \in \mathrm{N}_{\mathrm{G}}(\mathrm{T})\) ，\(w\) 是它在 W 中的类；则对 \(\lambda \in \mathrm{X}(\mathrm{T})\) ，有 \(\rho(n)(\tilde{\mathrm{V}}_{\lambda}(\mathrm{T})) = \tilde{\mathrm{V}}_{w\lambda}(\mathrm{T})\) 且 \(\dim \tilde{\mathrm{V}}_{w\lambda}(\mathrm{T}) = \dim \tilde{\mathrm{V}}_{\lambda}(\mathrm{T})\) 。

取 \(T' = T\) 、g = n 的公式 (9) 蕴含 \(\mathrm{P}(\rho, \mathrm{T})\) 在 w 下稳定；再者，\(\tilde{\rho}(n)\) 诱导一个从 \(\tilde{\mathrm{V}}_{\lambda}(\mathrm{T})\) 到 \(\tilde{\mathrm{V}}_{w\lambda}(\mathrm{T})\) 上的同构（公式 (8)），命题得证。

命题 6. 同态 \(\rho : \mathbf{G} \to \mathbf{GL}(\mathbf{V})\) 单当且仅当 \(\mathrm{P}(\rho, \mathrm{T})\) 生成 \(\mathbf{Z}\) -模 \(\mathrm{X}(\mathrm{T})\) 。

同态 \(\rho\) 单当且仅当它在 T 上的限制单（§2 第 6 节，命题 9）。再者，由于典则同态 \(\mathbf{GL}(\mathrm{V}) \to \mathbf{GL}(\tilde{\mathrm{V}})\) 单，我们可以把 \(\rho\) 换成 \(\tilde{\rho}\) 。于是由 (7)，\(\rho\) 在 T 上的限制的核是 \(\mathrm{P}(\rho, \mathrm{T})\) 中诸元素之核的交。从而结论由第 1 节命题 2 得出。

线性表示 \(\mathrm{L}(\rho)\) （t 在 \(\mathfrak{gl}(\tilde{\mathrm{V}})\) 中的）扩张成一个 C-李代数同态

\[
\tilde {\mathrm{L}} (\rho): \mathfrak {t} _ {\mathbf {C}} \to \mathfrak {g l} (\tilde {\mathrm{V}}).
\]

此外，回忆（第 1 节）我们对 X(T) 的每个元素 \(\lambda\) 联系了一个线性形式 \(\delta (\lambda)\) （\(\mathbf{t}_{\mathbf{C}}\) 上的），使得

\[
(\exp_ {\mathrm{T}} x) ^ {\lambda} = e ^ {\delta (\lambda) (x)}, \qquad x \in \mathfrak {t}.\tag{10}
\]

最后回忆（第 VII 章，§1，第 1 节）：对任一映射 \(\mu : \mathfrak{t}_{\mathbf{C}} \to \mathbf{C}\) ，我们用 \(\tilde{\mathrm{V}}_{\mu}(\mathfrak{t}_{\mathbf{C}})\) 表示 \(\tilde{\mathrm{V}}\) 中由诸 \(v\) （满足 \((\tilde{\mathrm{L}}(\rho)(u))(v) = \mu(u).v\) 对一切 \(u \in \mathfrak{t}_{\mathbf{C}}\) ）组成的向量子空间。

现在由 (7) 与同前命题 3 推出：

命题 7. a) 对一切 \(\lambda \in \mathrm{X(T)}\) ，有 \(\tilde{\mathrm{V}}_{\lambda}(\mathrm{T}) = \tilde{\mathrm{V}}_{\delta (\lambda)}(\mathbf{t}_{\mathbf{C}})\) 。

\begin{enumerate}
\def\labelenumi{\alph{enumi})}
\setcounter{enumi}{1}
\tightlist
\item
  映射 \(\delta : \mathrm{X(T)} \to \mathrm{Hom}_{\mathbf{C}}(\mathfrak{t}_{\mathbf{C}}, \mathbf{C})\) 诱导一个从 \(\mathrm{P}(\rho, \mathrm{T})\) 到 \(\mathfrak{t}_{\mathbf{C}}\) 在 \(\tilde{\mathrm{V}}\) 上的权所成的集上的双射。
\end{enumerate}

首先注意，若 W 通过把 W 的元素 w 对应为 \(t_{C}\) 的自同态 \(\mathrm{L}(w)_{(\mathbf{C})}\) 而作用在 \(t_{C}\) 上，则映射 \(\delta\) 与 W 在 X(T) 及 \(\mathrm{Hom}_{\mathbf{C}}(t_{\mathbf{C}}, \mathbf{C})\) 上的作用相容。

现在假定 \(k = \mathbf{R}\) 。用 \(\sigma\) 表示 \(\tilde{\mathrm{V}}\) 相对于 \(\mathrm{V}\) 的共轭，它由 \(\sigma(x + iy) = x - iy\) （\(x, y\) 属于 \(\mathrm{V}\) ）定义；对每个复向量子空间 \(\mathrm{E}\) （\(\tilde{\mathrm{V}}\) 的），\(\tilde{\mathrm{V}}\) 的在 \(\mathbf{R}\) 上有理且包含 \(\mathrm{E}\) 的最小子空间是 \(\mathrm{E} + \sigma(\mathrm{E})\) 。特别地，对一切 \(\lambda \in \mathrm{X}(\mathrm{T})\) ，存在实向量子空间 \(\mathrm{V}(\lambda)\) （\(\mathrm{V}\) 的），使得子空间 \(\mathrm{V}(\lambda)_{(\mathbf{C})}\) （\(\tilde{\mathrm{V}}\) 的）是 \(\tilde{\mathrm{V}}_{\lambda}(\mathrm{T}) + \tilde{\mathrm{V}}_{-\lambda}(\mathrm{T})\) （注意 \(\sigma(\tilde{\mathrm{V}}_{\lambda}(\mathrm{T})) = \tilde{\mathrm{V}}_{-\lambda}(\mathrm{T})\) ）。我们有 \(\mathrm{V}(\lambda) = \mathrm{V}(-\lambda)\) ，且诸 \(\mathrm{V}(\lambda)\) 是 \(\mathrm{T}\) 在 \(\mathrm{V}\) 上的、由 \(\rho\) 诱导的表示的同构型分支。

